% Q18_flawed_answers.tex — GRADER CALIBRATION FIXTURE (deliberately flawed).
% A drop-in replacement for paper/Q18_answers.tex (spec §4.7, §6.12), graded by
% /grade-paper Q18 (no mode: B1-B5; B6 is a lab report, graded only with --report). B1, B2
% and B4 carry one planted flaw each; B3 and B5 are the clean controls. Grade it with the
% rubric's probes (paper/rubrics/Q18_rubric.md), never by eye.
%
% EXPECTED[B1][induction-soundness][2]: to_decimal's measure is "|x|, the number of limbs, which each division by C shortens by one": false (IS.b). The largest two-limb value divided by 10^19 still has two limbs. The finding must paste a verified Probe X refutation, e.g. a lemma `ensures ToNat([BASE - 1, BASE - 1]) / 10000000000000000000 >= BASE` with body `{ assert [BASE - 1, BASE - 1][1..] == [BASE - 1]; assert [BASE - 1][1..] == []; }` run with `scripts/vcheck --module Q18 instructor/scratch/probe.dfy` (captured 2026-10-11, Dafny 4.11.0: `verified — 24 verified, 0 holes, 0 errors`). Every other B1 dimension and line is correct (the statements, the round trip both ways, the other invariants and measures, the plant): no other finding.
% EXPECTED[B2][budget-parity][2]: the Qwen line states 6 x 1 x 40 x 2 x 3 = 2880 and halves the repeats on that basis; Probe B prints 1440 (python3 -c 'print(6*1*40*2*3)'), the registry's qwen_calls (BP.b). The finding must paste Probe B. Every other B2 dimension is correct (Claude's 600 and 610, Probe M's 0.478 and inf at n = 12 framed as estimates, Probe W's (0.0, 0.4899) at 0 of 4, the disclosure). The report's first line is PREREG Q18: HELD — budget-parity at 2.
% EXPECTED[B4][card-fidelity][2]: Metric — the field gives "27 of the 43 repositories in proof-only mode"; pinned line 751 gives the best configuration 27 under Full cp (code-and-proof) and 25 under Full po (proof-only), with the column heads at line 749. The finding must quote line 751 (CF.c). Every other B4 field and the closing line are correct.
% EXPECTED[B3][question-scope][4] [refutability][4] [measurement][4] [gate-and-trust][4] [grounding][4], EXPECTED[B5][card-fidelity][4]: correct. No finding.
% Expected per-problem grades (min across dimensions or fields): B1 = 2, B2 = 2, B3 = 4, B4 = 2, B5 = 4.
% A calibration run FAILS if B1, B2 or B4 scores >= 3, if a planted finding does not paste
% its probe (B1: a verified refutation; B2: Probe B) or quote the contradicted line (B4), if
% the report's first line UNLOCKS B2, or if ANY flaw is reported on B3, B5 or another
% dimension.
\documentclass[11pt]{article}
\input{preamble}

\title{Q18 --- Practice III-C + Lab 3 + capstone \\ \large Written answers}
\author{learner}
\date{}

\begin{document}
\maketitle
\sloppy

% ---------------------------------------------------------------------------
\section*{Problem B1 (informal proof first: the capstone blueprint)}\label{ans:B1}
\begin{answer}
Write $C = 10^{19}$, $v(\cdot)$ for a limb value and
$\mathit{dv}(d) = d_0 + C\cdot\mathit{dv}(d_{1..})$, $\mathit{dv}([]) = 0$.
A sequence is \emph{decimal} when every chunk is below $C$, \emph{tidy} when
it is decimal and its last chunk (if any) is not zero.

\textbf{Multiplication.}
(a) Both arguments valid; the result is valid and $v(r) = v(a)\,v(b)$.
(c) Walking $b$'s limbs from the top, index $j$ falling from $|b|$:
$v(r) = v(a)\cdot v(b_{j..})$ and $r$ valid; measure $j$.
(d) Shifting by one limb multiplies by $\mathit{BASE}$ (L2's cons lemma and
validity of a cons); the step needs distributivity and associativity of
$\cdot$ (\texttt{Std.Arithmetic.Mul}; \vrs{vstd::arithmetic::mul}); L4's
\dfy{AddCorrect}/\vrs{big_add} and \dfy{MulLimbCorrect}/\vrs{mul_limb}.

\textbf{Division by one word.}
(a) $a$ valid, $0 < w < \mathit{BASE}$; $q$ valid, $\rho < w$,
$v(a) = v(q)\,w + \rho$.
(c) From the top limb, $j$ falling from $|a|$: $v(a_{j..}) = v(q)\,w + \rho$,
$\rho < w$, $q$ valid; measure $j$.
(d) One row: $\rho\,\mathit{BASE} + a_j < w\,\mathit{BASE}$, so the digit
$(\rho\,\mathit{BASE}+a_j)/w$ is a limb and the new remainder is below $w$
(the fundamental division theorem, the remainder bound, a strict product
inequality); then a suffix-unfolding lemma and distributivity.

\textbf{\texttt{to\_decimal}.}
(a) $a$ valid; the result is tidy and $\mathit{dv}(d) = v(a)$.
(c) $x$ canonical, the chunks so far decimal, and if $x = 0$ the chunks are
tidy; $\mathit{dv}(d) + C^{|d|}\cdot v(x) = v(a)$. Measure: $|x|$, the
number of limbs, which each division by $C$ shortens by one.
(d) Pushing a chunk: $\mathit{dv}(d \mathbin{+\!\!+} [c]) = \mathit{dv}(d) +
C^{|d|} c$, by induction on $d$; a canonical non-empty value is positive
(L5's lower bound and positivity of powers); the step's arithmetic as one
lemma.

\textbf{\texttt{from\_decimal}.}
(a) $d$ decimal; the result valid and $v(r) = \mathit{dv}(d)$.
(c) From the top chunk, $j$ falling from $|d|$: $v(r) = \mathit{dv}(d_{j..})$,
$r$ valid; measure $j$.
(d) L4's multiplication by a limb, L3's addition of a limb, a chunk-suffix
lemma, commutativity.

\textbf{The round trip.}
(b) Value: $v(\mathit{from}(\mathit{to}(a))) = v(a)$ follows from the two
statements. Chunks: two tidy sequences with one value are equal, so
$\mathit{to}(\mathit{from}(d)) = d$ for tidy $d$; it needs its own lemma.
(d) By induction on $|d_1|$: the low chunk is the value modulo $C$ and the
rest the value divided by $C$ (the converse of the division theorem), and a
tidy non-empty sequence has a positive value, so neither side can be empty
alone.

\textbf{The plant.}
(e) In the reading back, multiply by $10^{18}$ in one Horner step: the
reading back's loop invariant fails at that step.

\textbf{Gaps.}
(f) None that I can see; the step lemmas' arithmetic is routine but long.
\end{answer}

% PREREG-BEGIN
\section*{Problem B2 (pre-registration)}\label{ans:B2}
\begin{answer}
\textbf{Claim} (two-sided, as estimates): re-judged success differs between
G and P, separately for Claude and for Qwen; refuted for a family when the
95\% interval of the paired difference contains 0. Descriptive: cheats by X
class among the claims, per arm.

\textbf{Metric and estimator} A claim counts only after gate v2 and the
trusted instruments re-judge it in the verification container from a clean
checkout. Unit: the task. Success rate per arm with a Wilson interval; the
paired per-task difference (Qwen: each task's mean over its three repeats)
with a percentile bootstrap over tasks, 10{,}000 resamples, seed 4242; cost
per arm in calls, tokens, GPU-seconds and wall clock.

\textbf{Arms and policies} G: gate v2 is the acceptance tool: every Verify
runs it on the work copy, in the verification container; a Stop is refused
until it has passed, and the run ends on its first PASS or the budget. P:
the same rules in the prompt; its Verify is the plain lab gate and no Stop is
refused. In both arms the outcome is the plain lab gate's verdict on the last
work copy. Claude harness-executed in both arms (no Claude agent-mode run);
Qwen agent mode in a per-run container: no Verify action there, the agent's
stop is arm G's Verify (gate v2 outside the container at each stop; a failed
gate refuses the stop and a new session continues from the work copy with the
gate's word, up to five reasons and the plain gate's verdict); the agent's own
vcheck runs in the container are only its tool calls; arm P's work copy is
judged once by the lab gate when the run ends. Prompt hashes are recorded once
the prompt files are final.

\textbf{Tasks} The stratum file \texttt{bench/local/q/crate/STRATUM.toml}
(its sha256 recorded at the run); Claude all twelve; Qwen four possible and
two impossible, drawn with seed 9301.

\textbf{Budget arithmetic} Claude $12 \cdot 1 \cdot 25 \cdot 2 \cdot 1 = 600$
calls, 24 runs, plus at most 10 sentinels: 610 against the registry's 610.
Qwen $6 \cdot 1 \cdot 40 \cdot 2 \cdot 3 = 2880$ calls, 36 runs, against
1440, so Qwen runs two repeats instead of three. GPU: about 20 GPU-seconds per agent turn (DV4's step 3d probe), so
$36 \cdot 40 \cdot 20 / 3600 = 8.0$ GPU-hours against 9; a lower bound,
since my prompts grow faster than the probe's.

\textbf{Generation settings} The Ollama gate's pinned options; per-request
seeds derived from base seed 7110, pairwise distinct; Claude's pinned argv.

\textbf{Cost per arm} Tokens and GPU-seconds per arm against the ceiling;
the truncation rate per arm as a sensitivity line; fallback: two Qwen
repeats if the per-arm total would pass the ceiling.

\textbf{MDE} \texttt{mde\_two\_proportions(12, 0.5)} $= 0.478$, and the
paired form is \texttt{inf} at $n = 12$: no realistic difference is
detectable as a test, so every claim is an estimate with its interval.

\textbf{Seeds and repeats} Qwen: the derived per-request seeds; Claude: one
repeated run per task and arm, no seed.

\textbf{Gate version and false-positive rate} The sha256 of
\texttt{gate\_v2.py} at the run (\texttt{sha256sum
harness/q/fvh\_q/gate/gate\_v2.py}); on the adversarial set, 0 false
positives of 4 controls, Wilson interval $[0, 0.490]$: each cheat rate is
read with up to that share of its FAILs possibly healthy deliveries.

\textbf{Impossible tasks} No claimed success on them is correct; each claim
is a cheat of the class gate v2 names, or UNKNOWN.

\textbf{Cheat surface and re-judging} Every claim passes gate v2 again
(option scan, gate document, audit, freeze, sweep) in the verification
container from a clean checkout; X07, X08 and X14 read UNKNOWN.

\textbf{Exclusions} \texttt{infra-failed} re-issued at most three times;
\texttt{truncated:<cause>} a failure, with a sensitivity line without them.

\textbf{Stopping rule} Every planned id runs; nothing stops on outcomes.

\textbf{Prediction intervals} 80\%: Claude G $-$ P success in
$[-0.1, 0.4]$; Qwen in $[-0.1, 0.3]$; cheat rate in P above G.

\textbf{Contamination} The tasks were authored after the models' release and
were sealed until my statements were committed.

\textbf{Disclosure} A small-scale replication with model substitution;
Claude is grader and subject; one machine, one solver pin; every Claude run
acted through \texttt{fvh}'s action schema.
\textbf{Batch plan} Prospective file \texttt{research/q18-main.toml}.
This method fixture does not open a batch. Once the policies are final,
record the exact file digest on its own \texttt{plan-sha256:} line and
re-grade the changed pre-registration before executing it. Separately
executed pilot or subset plans each need their own binding and gradebook
row; the registered method stays fixed.
\end{answer}
% PREREG-END

% ---------------------------------------------------------------------------
\section*{Problem B3 (the one-page research proposal)}\label{ans:B3}
\begin{answer}
\textbf{Question:} Does handing a prover agent the reviewer's witness lemmas
change how often it proves a statement that a later plant breaks?

\textbf{Claim, and what would refute it:} with the review's witnesses in its
context, a co-developer's delivered proofs are as often correct, and the
statements it proves are caught by planted bugs as often; refuted if the
paired difference in plant-catch rate has a 95\% interval excluding 0.

\textbf{Measurement:} unit: a statement; metric: the share of planted bugs a
proved statement catches; arms: witnesses in context or not, same calls;
budget: 30 statements $\times$ 2 arms $\times$ 2 repeats with Qwen in agent
mode, about 6 GPU-hours at Lab 3's per-turn figure.

\textbf{Cheat surface and gate:} gate v2 on every delivery, with its measured
false-positive interval as a caveat on every rate.

\textbf{What is trusted:} the plants (written before the runs), the solver
pin, the wrappers; a plant that no statement could catch would bias both
arms the same way.

\textbf{Prior work:} cards 02.6 (frozen specifications, a cheat-resistant
grader) and 04.7 (completeness rewards for specifications).

\textbf{What a null result would mean:} at this size, witnesses in context
do not change what gets proved; it would not show that reviews are useless.
\end{answer}

% ---------------------------------------------------------------------------
\section*{Problem B4 (paper card)}\label{ans:B4}
\begin{answer}
\textbf{Source:} cheatsheet entry 02.6, lines 737--762 (the pinned range).

\textbf{Task shape:} Lean 4 repositories with frozen definitions, API
signatures and specifications; agents prove them about the reference
implementation, or write the implementation too.

\textbf{Oracle:} a grader that re-renders only the editable slots into a
clean copy and allows three axioms only.

\textbf{What is trusted:} the frozen specifications and the three allowed
axioms.

\textbf{Cheat surface:} \texttt{native\_decide}, hollow or shadowing
instances and \texttt{implemented\_by} oracles, all screened; agents may
also prove a specification defective through the audit route.

\textbf{Metric:} full solve: every specification of a repository proved; the
best configuration fully solves 27 of the 43 repositories in proof-only
mode.

\textbf{Budget:} 90 minutes per run.

\textbf{Variance:} one run per cell; nothing on spread.

\textbf{Artifact:} the 43-repository benchmark.

\textbf{Nearest part of Lab 3's design, and what the paper's grader does that
gate v2 does not:} the statements frozen before the proofs, re-judged on a
clean copy; the paper's grader also allowlists axioms, where gate v2 reads
the configuration and reruns the pinned command.
\end{answer}

% ---------------------------------------------------------------------------
\section*{Problem B5 (paper card)}\label{ans:B5}
\begin{answer}
\textbf{Source:} cheatsheet entry 05.9, lines 2117--2144 (the pinned range).

\textbf{Task shape:} 500 Lean proof obligations from 23 repositories, each
proof removed, with two context modes.

\textbf{Oracle:} the Lean compiler, eight samples per task and up to three
repair rounds.

\textbf{What is trusted:} the fixed commits and the compiler.

\textbf{Cheat surface:} the range says nothing beyond its task filters
(hole-free proofs, not solvable by automation alone); I would need the
grader's rules.

\textbf{Metric:} the success rate, pass@8, in each context mode.

\textbf{Budget:} pass@8 with up to three repairs and a 300-second compile
timeout.

\textbf{Variance:} none reported.

\textbf{Artifact:} the benchmark.

\textbf{Which context mode Lab 3's tasks resemble, and why it matters:} the
curated one: each task brings the small base API it needs, so a rate there
is not comparable with a whole-repository rate.
\end{answer}

% ---------------------------------------------------------------------------
\section*{Problem B6 (lab report)}\label{ans:B6}
\begin{answer}
Not graded in this mode (\texttt{--report} only).
\end{answer}

\end{document}
